\chapter*{Resumen}
\addcontentsline{toc}{chapter}{Resumen}
Esta tesis propone un marco metodológico para el sensado humano mediante infraestructura inalámbrica existente. Una medición de canal no proporciona una observación directa de la presencia, la postura o la fisiología, sino el resultado de una cadena constituida por la geometría, la propagación electromagnética, la forma de onda, la instrumentación, el protocolo, el ruido y la interferencia. El trabajo plantea representar dicha cadena mediante un gemelo digital inalámbrico diferenciable y físicamente estructurado, calibrarlo con mediciones reales e invertirlo con una cuantificación explícita de la incertidumbre y la observabilidad. Wi-Fi y LoRa se estudian como observadores complementarios de un mismo estado físico. La validación progresa desde casos analíticos y un gemelo estático hasta el movimiento controlado, la respiración, la presencia de varios usuarios, la transferencia de la simulación a condiciones reales y el diseño de redes distribuidas. A lo largo del manuscrito se distinguen las contribuciones propuestas de los resultados que aún deben demostrarse experimentalmente.

Los experimentos documentados abarcan S1/B0--B2 y S4.1--S4.6d.2. Se utiliza un banco de 10\,000 posiciones de DICHASUS dc01, dividido inicialmente en 5000 posiciones de entrenamiento, 100 de validación y 4900 de prueba. Los materiales neuronales B2 reducen el MAE de potencia de 6.29 a 2.22 dB y el de dispersión RMS de retardos de 43.21 a 16.50 ns respecto de B0, sin mejorar la fase bruta. La descomposición posterior distingue una pendiente espectral efectiva, un desfase constante por enlace y un residual con estructura espectral. Sus asociaciones con descriptores energéticos se examinan mediante análisis entre posiciones, dentro de posición y remuestreo agrupado. La predicción del retardo se evalúa primero bajo interpolación densa y después mediante cinco bloques espaciales con exclusión de 0.5 m. En esta última evaluación, el predictor que combina posición y descriptores del trazado alcanza un MAE medio de retardo de 18.69 ns, frente a 27.56 ns al usar sólo posición. Aplicado al canal, y manteniendo ajuste del CPO con la medición evaluada, obtiene una mediana de error angular de 64.08 grados y coherencia de 0.448. S4.6 caracteriza 640\,000 fases residuales y recupera el tiempo de 10\,000 capturas. La resultante global es 0.001088, mientras que la concentración mediana dentro del arreglo es 0.4216. Las predicciones por vecinos espaciales y temporales de la dirección común producen errores medianos próximos a 90 grados, sin ventaja clara frente al control aleatorio. Estos resultados apoyan una distinción operacional entre retardo parcialmente predecible y fase común que requiere un tratamiento de referencia explícito; no identifican su causa física ni demuestran impredecibilidad universal. La evaluación de invariantes espaciales y espectrales S4.7a, la incertidumbre por trayectoria S4.7b, las representaciones alternativas y la fidelidad diferencial constituyen los siguientes bloques metodológicos.

La propuesta se articula ahora mediante dos ramas complementarias: una robusta frente
a transformaciones instrumentales identificadas y otra coherente que preserve respuesta
al movimiento. Se propone contrastar geometría manual, RGB-D, LiDAR, regularización RF
y calibración mediante G0--G4, separando precisión global, registro de antenas y referencia
local de desplazamiento. La campaña E0--E5 progresará desde estabilidad instrumental,
reflector mecánico y fantoma respiratorio hasta personas y transferencia. Referencia,
repetibilidad y sensibilidad se validarán antes de interpretar el problema inverso humano.
Estos bloques actualizan la metodología y la propuesta de trabajo; sus resultados aún
no forman parte de la evidencia experimental documentada.
